%!TEX root = parita-msc.tex

\chapter{Conclusion and Future Work}
\label{chap:conclusion}

\begin{doublespace}

%This Chapter focuses on the prime offerings of this research work where one of the sections includes the conclusions drawn from the work done and the other discusses future work that can be done to succeed in this research.

This Chapter presents a summary of the research undertaken and 
contained in this document.
It identifies the contributions made in this thesis and opportunities for future work.

\section{Conclusion}

After motivating the topic in Chapter 1,
Chapter 2 provides a detailed background study on the fundamental topics
supporting \ac{lod} and the tools which facilitate its publication on the web.
These include 
Semantic Web, 
Linked Open Data, 
\ac{rdf}, and 
\ac{owl} ontologies. 
This thesis describes the lifecycle of transforming data contained in a spreadsheet
into \ac{lod} that is available on the web of data, 
both in general terms (Chapter 3) 
and in the specific case of a student feedback survey (Chapter 4).
%a new attempt on developing a linked open data instrument for feedback surveys. 
%It focuses on the issues encountered and the user experience of various tools and references used for the modelling of the instrument and for data visualization. 
%The thesis revolves around some fundamental topics like 
%Semantic Web, 
%Linked Open Data, 
%\ac{rdf}, and 
%\ac{owl} ontologies. 
%It also describes how and why Protégé has been used over other ontology development tools to describe survey vocabularies. 
The thesis can be considered as a guide for researchers who are interested in 
transforming their own data into \ac{lod}.
%exploring this lesser known aspect of semantic web.

%the usability of protégé, and 
%discussion on some tools for linked open data visualizations. 
%Chapter~\ref{chap:methodology} describes the methodology followed to 
%describe, convert, visualize, and publish linked open data for survey in Chapter~\ref{chap:proofOfConcept}
%using protégé, OpenRefine, and the visual representation of the ontology 
%toward the end.

%thesis reviews the literature around publishing 

\section{Future Work}

Modelling of the range of Likert~\cite{what-is-Likert} scales.
with different labels and values.

Jekyll-RDF integration.

Conversion and management of all feedback data collected by Hepting

https://github.com/schimatos/schimatos.org -- data entry
Developing a data entry form for the feedback survey instrument
that can be defined in part ...


%This thesis also offers a number of future development options like building a real-time linked open data application for feedback surveys that 
%follows the linked data life cycle. 
%
%Web Scraping (or Web Extraction) 
%is a method of extracting desired data from the resources on the \ac{www} 
%into a separate file for data retrieval and data analysis. 
%Data Extraction is one of the steps in the linked data life cycle that can be explored in future research to facilitate use of data that has already been published to the web with \ac{html} tables.

Versioning of the survey ontology is a topic of interest. 
The ontology vocabularies can be reused for appropriate applications but they may need to be updated over time. 
Hence, 
ontology versioning and version management of the survey instrument for this feedback survey, 
as well as for other types of surveys is an aspect of this research that can be explored.

directly from the survey ontology would facilitate recording of survey responses collected offline.
\end{doublespace}